\section{Related Work}
\label{sec:related}

\subsection{Image generation and structural control}

Diffusion models learn to reverse a gradual noising process and provide a stable
foundation for high-quality image synthesis \citep{ho2020ddpm}. Improved
parameterizations and sampling methods made the approach more efficient
\citep{nichol2021improved}, while latent diffusion reduced the cost of
high-resolution synthesis by operating in a learned latent space
\citep{rombach2022latent}. Joint image-text representations such as CLIP enabled
general semantic comparison between prompts and images \citep{radford2021clip}.
These developments explain why an existing image model can serve as a strong
source-art generator in \skinsmith.

The limitation is spatial control. A text prompt can describe a dragon, cloud, or
metallic surface, but it does not specify how those elements should survive a
particular UV atlas. ControlNet demonstrates the broader value of supplementing
text with spatial conditions such as edges, depth, and segmentation
\citep{zhang2023controlnet}. \skinsmith\ adopts the same motivation without
training a new conditioning network: semantic planning is separated from
deterministic geometry-based placement and evaluation.

\subsection{Texture synthesis, periodicity, and UV atlases}

Classical texture-synthesis methods grow or assemble images from local
neighborhoods and patches \citep{efros1999texture,efros2001quilting}. Graph-cut
textures improve joins between overlapping patches by optimizing the cut through
the overlap \citep{kwatra2003graphcut}. Periodic-plus-smooth decomposition
addresses discontinuities caused when a finite image is treated as periodic in
the Fourier domain \citep{moisan2011periodic}. This work supports the use of a
tileable pattern as a meaningful baseline and motivates the frequency-domain
border repair used by \routeA.

Square-border periodicity is nevertheless different from continuity on a mesh.
Surface parameterization flattens a three-dimensional surface into one or more
charts and must trade distortion against chart layout
\citep{floater2005parameterization}. Least-squares conformal maps are a widely
used approach to low-angle-distortion atlases, but chart boundaries still
interrupt large regular patterns \citep{levy2002lscm}. TextureMontage further
shows that positioning multiple images on an arbitrary surface requires explicit
control over distortion and transitions \citep{zhou2005texturemontage}. In a game
pipeline, the UV contract is often already fixed. \skinsmith\ therefore does not
replace the parameterization; it derives geometry maps and seam correspondences
from the bound asset and works within that atlas.

\subsection{Diffusion-based 3D texturing}

Several systems use two-dimensional diffusion priors to generate appearance for
three-dimensional objects. DreamFusion demonstrated that rendered views can
connect a text-to-image prior to a 3D representation without a large paired 3D
dataset \citep{poole2022dreamfusion}. TEXTure progressively paints an existing
mesh from depth-conditioned views and supports text-guided generation and editing
\citep{richardson2023texture}. Text2Tex similarly uses a progressive
generate-and-refine process with view partitioning and automatic view selection
\citep{chen2023text2tex}. TexFusion aggregates denoising predictions in a shared
latent texture representation to improve global consistency
\citep{cao2023texfusion}. Point-UV Diffusion combines a low-frequency point-space
representation with UV-space refinement \citep{yu2023pointuv}. Paint3D adds UV
inpainting and high-resolution refinement to produce complete, lighting-less
texture maps \citep{zeng2024paint3d}. Synchronized multi-view diffusion shares
denoised content between overlapping views during sampling
\citep{liu2024syncmvd}.

This literature establishes that UV discontinuity, view inconsistency,
incomplete coverage, and baked illumination are central problems. It also limits
the novelty claim that can reasonably be made here: \skinsmith\ is not the first
system to texture a mesh using diffusion. Its focus is a different layer of the
problem---an asset-bound, human-supervised Agent that begins with compact intent,
retains the target game's UV contract, exposes mapped alternatives before formal
execution, enforces deployment constraints, and preserves the evidence behind
every decision.

\begin{table}[t]
  \centering
  \caption{Positioning against the closest 3D-texturing systems. ``Not central''
  means that the property is not a principal reported contribution, not that an
  extension would be impossible.}
  \label{tab:closest-work}
  \small
  \begin{tabularx}{\textwidth}{p{1.8cm} X X X}
    \toprule
    Work & Main mechanism & Consistency and UV treatment & Human and deployment control\\
    \midrule
    TEXTure & Iterative depth-conditioned painting & Successive views are projected
    to a mesh texture & Text-guided generation and editing; hard deployment gates
    are not central\\
    Text2Tex & Progressive generation and refinement & View partitioning and
    automatic next-view selection update visible texels & Prompt-controlled; no
    explicit rollback contract reported\\
    TexFusion & Shared latent texture denoising & View predictions are aggregated
    on one latent texture & Text-conditioned; no staged human checkpoint contract\\
    Point-UV & Point-space coarse colour followed by UV diffusion & Low-frequency
    global condition plus UV refinement & Conditional generation; no game-specific
    export gate\\
    Paint3D & Multi-view fusion, UV inpainting, and UVHD refinement & Complete
    high-resolution lighting-less UV maps & Text or image conditioning; focuses on
    texture quality rather than Agent rollback\\
    SyncMVD & Synchronized denoising across overlapping views & Shares latent
    content at each denoising step & Text-guided; no three-stage selection contract\\
    \skinsmith & User-selected master artwork and deterministic weapon-space bake &
    Preserves the authoritative UV; evaluates fixed mapped views and true seams &
    Three checkpoints, hard feasibility, one bounded correction, and rollback\\
    \bottomrule
  \end{tabularx}
\end{table}

\subsection{Rendering and evaluation}

Differentiable rendering research formalizes the link between a 3D representation
and image-space supervision \citep{kato2018renderer,liu2019softras}. \skinsmith\
does not optimize through rendering gradients, but follows the same general
principle: the mapped views are evidence about the texture that cannot be recovered
from the source image alone.

Perceptual metrics also require care. LPIPS showed that learned feature distances
often correlate with human perceptual judgments better than simple pixel
distances \citep{zhang2018lpips}. CLIPScore provides reference-free semantic
compatibility, while documenting cases in which such a score is weaker
\citep{hessel2021clipscore}. Accordingly, \skinsmith\ keeps texture, seam,
multi-view, and semantic evidence separate. Automated review recommends and
explains; it does not replace artwork selection by the user.

\subsection{Constrained selection and bounded Agents}

Weighted aggregation is convenient, but a high score on one objective can hide a
violation on another. Feasibility-oriented constraint handling instead ranks
valid solutions ahead of invalid ones \citep{deb2000constraints}. NSGA-II
formalizes non-dominated sorting for multi-objective optimization
\citep{deb2002nsgaii}. \skinsmith\ adapts these principles to a small, expensive
candidate pool rather than running an evolutionary search: it filters hard
violations, computes the feasible Pareto set, and uses a fixed objective order.

Agent research provides a second relevant foundation. ReAct interleaves reasoning
and actions so observations can change the trajectory \citep{yao2023react}.
Self-Refine and Reflexion use feedback and memory to improve later attempts without
updating model weights \citep{madaan2023selfrefine,shinn2023reflexion}. These
methods motivate explicit observation, diagnosis, action, and memory, but they do
not imply that every revision is better. \skinsmith\ therefore limits refinement
to one local round and requires measured improvement before acceptance.

\subsection{Human control and responsible generation}

Human-AI interaction guidelines emphasize communicating capabilities, allowing
correction, and preserving user control when model behavior is uncertain
\citep{amershi2019guidelines}. Co-creative systems treat generation as a turn-based
collaboration \citep{guzdial2019cocreative}, while PromptChainer demonstrates the
value of visible intermediate stages for complex AI tasks
\citep{wu2022promptchainer}. These ideas inform the three checkpoints in
\skinsmith: theme, direction, and mapped artwork.

Finally, datasheets and model cards show why provenance, intended use, evaluation
conditions, and limitations should be recorded
\citep{gebru2021datasheets,mitchell2019modelcards}. Responsible-AI scholarship
also warns that fluent outputs can obscure data, bias, environmental, and
accountability concerns \citep{bender2021parrots}. The project consequently logs
model and provider identifiers, redacts API keys, preserves formal runs, prohibits
copied skins and protected characters, and keeps publication outside the tested
scope.

\subsection{Research gap}

Prior work provides strong components: text-to-image generation, spatial
conditioning, periodic texture synthesis, 3D-aware texture generation,
multi-objective optimization, and iterative feedback. The gap examined here is
their practical integration around an existing game asset:

\begin{quote}
How can compact user intent be converted into a human-selected source design that
is deterministically bound to an authoritative UV contract, evaluated after real
geometry mapping, and refined only within explicit feasibility and rollback rules?
\end{quote}

This framing treats deployability, evidence, and human authority as part of the
method rather than as post-processing around a generated image.
